\section{The Angel of Death Spread his Wings}

\begin{quotex}
\begin{verse}
For the Angel of Death spread his wings on the blast,\\ And breathed in the face of the foe as he passed:\\ And the eyes of the sleepers waxed deadly and chill,\\ And their hearts but once heaved, and for ever grew still!
\end{verse}
\flright{\textsc{Lord Byron}, \emph{The Destruction of Sennacherib}}
\end{quotex}


I'd like to follow up the discussion on Hope, and why it is a virtue rather than an emotion. Hope has two parts:

\begin{itemize}


\item There is a desire for something

\item There is the expectation of its being fulfilled
\end{itemize}


The expectation distinguishes Hope from wishful thinking or blind faith. What justifies that expectation? In a natural sense, there may be knowledge of cause and effect, or knowledge of conditions, that make hope reasonable. Some things may be under the control of your will, so you can legitimately expect to achieve your goal. In the latter case, the alternatives are not at all helpful. If you lack hope, you will not even try to achieve your goal. Or else wishful thinking might convince you that effort is unnecessary.

For supernatural Hope, the expectation is based on Providence or God's Will. That is of a different order, because the very means of fulfillment are not visible. The situation seemed hopeless to the Hebrews when King Sennacherib sent the Assyrian army laid siege to Jerusalem. Hezekiah prayed to God for help. He sent the angel of death who put to death 185,000 Assyrians as they slept.

I suppose it is a test of Hope to expect that God will smite his enemies.

\paragraph{Alternatives to Hope}


Now Hope is a virtue, since it is an act of will: one chooses to live in expectation of the fulfilled desire.

Cologero asked me about a situation in which there is a desire accompanied by the fear of its being fulfilled. I don't think there is a word for it. He suggested \emph{Sehnsucht}, the German word for ``longing''. I know that feeling, but I don't think his analysis is correct. It is as if the desire itself is unclear, that is, you know it but don't know it, like that word that is always on the tip of your tongue. The fear, if that is not too strong a word, arises from not knowing exactly what to expect.

\paragraph{Arabian Nights and Space Odyssey}


Cologero keeps delaying the march to 1001. He insists that he is dealing with a personal situation that suddenly arose and has diverted his attention. He even asked me today about the symbolic meaning of various flowers, as if I would know.

I suspect he believes Gornahoor is related in some way to his own mortality. We've been looking into alternatives, one of which would be to continue Gornahoor with a different focus and design. I suggested that Phase 1 would be called \textbf{Arabian Nights} and the next series of posts to \#2001 would be \textbf{Space Odyssey}.

His spiritual father, Losang Shenphen, has already approved.

\paragraph{The Meaning of a Library}


Cologero asked me to post a few notices. First of all, a library is not a mere heap of books. There is a coherence to a personal library. Some texts belong together, other texts reinforce each other. They reveal a history of influences over many years. Some are rare, even in our age of high-speed scanning and print-on-demand technologies. The point is that a library is a ``seamless garment'', all or nothing.

Several people have asked for individual volumes. How can he assign a few thousand books to different people. You guys need to think. Moreover, some of you have appeared a little too eager. Cologero has asked me to keep you out of his hospital room, if it ever comes to that, for fear you will cut off his life support.

The goal was always to give an independent scholar or a new professor a jump start for his own library, assuming, of course, his intellectual interests were in alignment.

\paragraph{Writer's Workshop}


In an effort to encourage more writers on esoteric and Hermetic topics, Cologero plans to hold a Writer's Workshop sometime next month. Some suggested topics will include:

\begin{itemize}


\item Idea density and specificity

\item How to write on different levels simultaneously

\item How to writing esoterically

\item Choosing a theme

\item Choosing a title

\item Promotion
\end{itemize}


Update 24 April 2021: Or probably not.


\flrightit{Posted on 2018-04-24 by Aeneas}

\begin{center}* * *\end{center}

\begin{footnotesize}
\begin{sffamily}
\texttt{TeenyTinyTarot on 2018-04-27 at 13:14 said:} 

``Cologero keeps delaying the march to 1001. . . . I suspect he believes Gornahoor is related in some way to his own mortality. We've been looking into alternatives... '' One possible alternative might be to review all your posts to see if there were some that might be deleted and/or combined with others in a way that would legitimately postpone the end while at the same time ``combing, curling, and braiding'' the remaining artifact for posterity, eh?

``Even at the age of 80, Plato did not stop combing, curling and braiding his dialogues. Every philosopher knows about the loving care he took with them.''  Dionysius of Halicarnassus

p.s. Apropos of hope, think ``confidence in the face of the unknown''. :-)

\end{sffamily}
\end{footnotesize}

